\chapter{Introduction}
This chapter introduces the motivation, research goal, and outline of this thesis. \Autoref{intro:motiv} motivates the research problem, \autoref{intro:goal} defines the specific research questions addressed, and \autoref{intro:outline} outlines the remainder of this work.

\section{Motivation}\label{intro:motiv}
\subsection{Two-pore domain potassium ion channel}
Cellular electrical excitability is tightly regulated by precise control of ion fluxes across the plasma membrane. Background or leak potassium currents play an essential role in stabilizing the resting membrane potential and opposing over-excitation in the nervous, cardiovascular, and immune systems \parencite{nichollsNeuronBrainCellular1992}. While systemic potassium imbalances such as hypokalemia or hyperkalemia cause severe clinical pathologies, genetic-mutational channelopathies of potassium channels are directly linked to various neurological and physiological disorders \parencite{turkaydinEnergeticCrosstalkFilter2026, baukalCalciumdependentActivationTREK12025, kimUnravelingMolecularReason2024}. In this context, potassium channels emerge as potential therapeutic targets and have been actively investigated. 
Among the diverse families of potassium channels, the two-pore domain potassium (\KtwoP) channels are prominent mediators of these background currents. The human genome encodes 15 \KtwoP\ subunits, which are classified into six distinct subfamilies: TWIK, TREK, TALK, TASK, THIK, and TRESK.

\section{Research Goal}\label{intro:goal}
Our goal in this study is to examine which point mutation can lead to the up or down conformation of hTREK-2. The experimental and computational data obtained from \textcite{musinszkiUpstateTransitionDefault2026} are used to analyze our TREK mutation assays, comparing ProteinMPNN and ESM-IF1. 

\section{Outline of Work}\label{intro:outline}
In this study, we investigated TREK-2, a member of the mechanosensitive TREK \KtwoP\ family. Previous structural studies have revealed that TREK channels undergo vertical displacements of the M4 transmembrane helix, which define two distinct up and down conformations, as illustrated in \autoref{fig:trek_updown} \parencite{nataleStructuralInsightsMechanisms2021, schmidpeterMembranePhospholipidsControl2023, turkaydinAtomisticMechanismCoupling2024}. This structural conversion alters the physical interface between the M2, M3, and M4 helices. During the transition to the up-state, the M2 and M3 helices of the same chain undergo buckling toward M4. Concurrently, M4 shifts toward and stabilizes against the M2' helix of the adjacent chain. Within this region, the P2 helix interacts with residue S1, whose correctly flipped backbone carbonyl increases channel currents. These transmembrane helix couplings are known to allosterically regulate ion gating in TREK channel proteins \parencite{proksNorfluoxetineInhibitsTREK22021, turkaydinAtomisticMechanismCoupling2024}.

Using a similar mechanism, one known TREK inhibitor, Norfluoxetine (NFx), binds within a known fenestration site (PDB ID: 4XDK) and inhibits ion gating by hindering M4 movement to M2'.  The mutations Y315A \parencite{proksNorfluoxetineInhibitsTREK22021} and L320W \parencite{dongK2PChannelGating2015}, located in the M4 helix, have shown that the allosteric inhibitory effect of NFx could be reduced. These mutations impair the allosteric movement of TM helices in filter gating, resulting in NFx insensitivity. Following this mechanism, it is interesting to investigate TREK mutations in this fenestration site that induce a permanent conformational change leading to anormal current behavior of TREK channels. 

\Textcite{musinszkiUpstateTransitionDefault2026} investigated human TREK-2 (hTREK-2) channel mutations using molecular dynamics (MD) simulations and measured NFx inhibition across complementary mutations in rat TREK-1. They used the up (PDB ID: 4BW5) and down (PDB ID: 4XDJ) crystal structures for the MD simulations and evaluated them by comparing conformational free energy changes $\Delta\Delta G$ between WT and mutants. Additionally, they quantified the relative open probability \textit{P$_o$} by comparing currents at pH 5 and pH 8. Mutants with high \textit{P$_o$} were subsequently measured for the effective inhibition concentration $IC_{50}$ of NFx and correlated with the calculated $\Delta\Delta G$ of hTREK-2. We integrate these experimental and computational datasets as a baseline to analyze our TREK mutation assays using the following method.